\section{The DIA Pre-Registration Defense}
\label{sec:preregistration}

\subsection{Why 0.3 Seats Is Not Enough Justification for No Defense}

Even though 0.3 seats of parameter-based variation is below the noise floor,
the DIA should still eliminate this gaming vector.
The reason is not the magnitude of the effect; it is the \emph{appearance}
of manipulation.
If parameters are chosen after Census data are released, an opponent can
always argue---even without evidence---that the parameter choices were
motivated by their partisan implications.
This argument is difficult to rebut in litigation: the defendant cannot
prove they did not observe partisan implications before setting parameters.

Pre-registration converts this he-said/she-said dispute into a
cryptographic fact.

\subsection{The Pre-Registration Mechanism}

The DIA's pre-registration protocol for parameters operates as follows:

\begin{enumerate}
\item \textbf{Twelve months before Census day}: The redistricting
  administrator files a parameter configuration document containing
  all five parameters and their values.
  The document is hashed with SHA-256 and the hash is filed with the
  Secretary of State.
  The hash is computed as:
  \[
    h = \mathrm{SHA256}(\texttt{ufactor} \| \texttt{acounty} \| T
    \| \texttt{aswing} \| \texttt{wvra} \| \texttt{algorithm\_binary\_hash}).
  \]
\item \textbf{Six months before Census day}: A public comment period
  opens.
  Any party can challenge parameter values as inconsistent with the
  statutory specification or as technically unjustified.
  The administrator must respond to all challenges.
\item \textbf{Census day}: Redistricting data released.
  Parameters cannot be changed.
\item \textbf{Computation day}: The algorithm runs using the pre-registered
  parameters.
  The manifest records the pre-registration hash alongside the computation
  output hash.
\item \textbf{Verification}: Any party can verify that the deployed
  parameters match the pre-registered hash using \texttt{bisect label-verify}.
\end{enumerate}

\subsection{What Pre-Registration Guarantees}

Pre-registration provides three guarantees:

\begin{theorem}[Pre-Registration Eliminates V1 Gaming]
Under the DIA pre-registration protocol, no parameter gaming vector
can produce a partisan outcome differing from the pre-registered
default by more than the seed-variance noise floor ($\sigma \approx 0.8$ seats).
\end{theorem}

\begin{proof}
The parameter values are fixed before Census data are released.
Any change to parameters after release changes the manifest hash.
The maximum effect of parameter variation within admissible ranges is
0.3 seats (from the 50-state sweep).
Since 0.3 $<$ 0.8 = $\sigma$, even the maximum effect is within
the noise floor.
$\square$
\end{proof}

Guarantee 1: No post-data parameter optimization is possible.
Guarantee 2: Any deviation from pre-registered parameters is
cryptographically detectable.
Guarantee 3: The maximum undetectable effect (from parameters within
the pre-registered values) is zero by definition.

\subsection{Challenge-Period Safeguard}

The six-month challenge period serves a second purpose beyond public
accountability: it recruits opponents as algorithm auditors.
Any party who suspects the pre-registered parameters are partisan-motivated
can challenge them during the comment period.
If the challenge identifies a credible partisan motivation (e.g., the
submitted $\acounty = 10$ is the value that minimizes Democratic seats
in a specific state), the administrator must either justify the choice
on technical grounds or resubmit different parameters.

This self-enforcing mechanism does not require a neutral arbitrator
to evaluate each parameter choice; it relies on adversarial scrutiny
to surface problems before they affect the map.
